\section{Field-valued points and coherent trace}\label{sec:coherent-trace}

Continue with $p\geq5$, $t=p-1$, $h=p$, and the pure subgroup $H$
of Theorem~\ref{thm:lower-hm}. To transport the finite upper layers,
we need exponent one for modules over the mixed algebra
\[
 D=E_t\wedge E_h,\qquad D'=L_{K(t)}D.
\]
Here $H$ acts on the first factor. Section~\ref{sec:exponent-transport}
will make each transformed layer a $D'$-module in $\mathcal D_H$.
The \(K(t)\)-acyclic spectra form a smash ideal. Monoidal localization
of \(D\) into \(\mathcal C_t\), followed by the lax monoidal inclusion
into ordinary spectra, therefore gives \(D'=L_{K(t)}D\) a commutative
algebra structure in \(\mathcal D_H\). Its localization map is
equivariant and is a map of these commutative algebras.

\begin{proposition}[The completed mixed trace]\label{wup:completed-trace}
The mixed $H$-algebra $D'$ has an underlying degree-zero class whose
$H$-trace is a unit.
Every \(D'\)-module in \(\mathcal D_H\) therefore has
\(H\)-exponent at most one.
\end{proposition}

A trace unit gives the required coherent retraction by induction and
coinduction. We first prove that implication. We then construct the
trace unit from the two formal-group isomorphisms supplied by the
factors of $D$: the fixed $E_h$ factor forces $H$ to act freely on
field-valued points of
\[
 R=\pi_0D/\m_t\pi_0D.
\]
The ideal is $H$-stable. The mixed cooperation calculation
of Lemma~\ref{wup:mixed-flat} makes $D$ even periodic, with an
orientation and periodic unit from $E_t$.

\begin{lemma}[Coherent trace splitting]\label{wup:trace-split}
Let \(D'\) be a commutative algebra in \(\mathcal D_H\). Suppose an
underlying class \(z\in\pi_0D'\) has unit trace
\[
 v=\sum_{g\in H}g(z)\in(\pi_0D')^\times.
\]
Then every \(D'\)-module \(U\) in \(\mathcal D_H\) has
\(\expn_H(U)\leq1\).
\end{lemma}

\begin{proof}
Since \(H\) is finite, the underlying object \(A_H\wedge U\) is
\(F(H_+,U)\), with diagonal action.
Untwisting the value at \(g\) by the action of \(g^{-1}\) on \(U\)
identifies it with the coinduction of the underlying spectrum:
\[
 \theta:A_H\wedge U\xrightarrow{\ \simeq\ }
       \operatorname{Coind}_{1}^{H}(\operatorname{Res}_{1}^{H}U).
\]
Indeed, on functions the diagonal action is
\((h f)(g)=h f(h^{-1}g)\), and \((\theta f)(g)=g^{-1}f(g)\)
turns it into translation \((h\theta f)(g)=(\theta f)(h^{-1}g)\).
For finite \(H\), the finite coproduct-product equivalence in spectra
identifies induction with coinduction. Induction adjunction takes
the underlying multiplication map \(\mu_z:U\to U\) to a map with
coherent \(H\)-action
\[
 r_z:A_H\wedge U\longrightarrow U.
\]

Let \(j:U\to A_H\wedge U\) be the unit map. Under \(\theta\), its
component at \(g\) is \(g^{-1}:U\to U\). The induction map has
component \(g\mu_z\) at \(g\), so the underlying composite is
\[
 r_zj=\sum_{g\in H}g\mu_zg^{-1}
     =\mu_{\sum_g g(z)}=\mu_v.
\]
The middle equality uses the compatibility of the \(D'\)-module
structure with the \(H\)-action. Since \(v\) is a unit, this is an
equivalence of underlying spectra. Equivalences of \(H\)-diagrams
are detected on underlying objects, so \(r_zj\) is an equivalence
in \(\mathcal D_H\). Hence \((r_zj)^{-1}r_z\) retracts \(j\)
there. Its fiber map is null, which is exponent at most one.
The class \(z\) was only an underlying class; no homotopy-fixed
lift of it is needed.
\end{proof}

\begin{lemma}[Freeness on field-valued points]\label{wup:field-points}
For every unital ring map \(x:R\to\Omega\) to a field, the equality
\(x\circ g=x\) for \(g\in H\) implies \(g=1\).
\end{lemma}

\begin{proof}
We first recall explicitly the naturality of the formal group used
here. For a complex-oriented even periodic commutative ring
spectrum \(B\), an orientation and a periodic unit give a coordinate
in degree zero. The finite projective-space calculation and its
surjective restriction tower give
\[
 B^0(\mathbb{CP}^{\infty})=(\pi_0B)[[z_B]].
\]
Products of projective spaces give the corresponding group law.
A unital ring-spectrum map \(f:B\to C\) carries an orientation
to an orientation and a periodic unit to a unit. Its coordinate
therefore has an invertible linear coefficient in the chosen
coordinate of \(C\). The inverse of this coordinate substitution
gives an isomorphism from the base change of the formal group of
\(B\) to the formal group of \(C\).
The product calculation gives compatibility with group laws,
and this construction respects composites of ring-spectrum maps.
All later base changes are base changes of these algebraic formal
groups along maps of coefficient rings.

Let \(\rho:\pi_0D\to\Omega\) be the composite through \(x\).
Write \(\iota_t:E_t\to D\) and \(\iota_h:E_h\to D\) for the
two factor maps, and let \(\rho_t,\rho_h\) be their coefficient
composites with \(\rho\). The map \(\rho_t\) factors through
\[
 O_t/\m_t=k_t\longrightarrow\Omega.
\]
The last map is injective, since it is a unital map from a field.
After base change along \(\rho_t,\rho_h,\rho\), the factor maps
give formal-group isomorphisms
\[
 \phi_t:\mathcal F_{t,\Omega}\xrightarrow{\ \cong\ }
              \mathcal F_{D,\Omega},\qquad
 \phi_h:\mathcal F_{h,\Omega}\xrightarrow{\ \cong\ }
              \mathcal F_{D,\Omega}.
\]

Let \(g_D\) denote the action of \(g\) on \(D\).
The equality \(x\circ g=x\) implies
\(\rho\circ\pi_0(g_D)=\rho\), so naturality makes \(g_D\) an
automorphism \(u_g\) of the single formal group
\(\mathcal F_{D,\Omega}\). Since \(g_D\iota_h=\iota_h\),
naturality with the second factor gives \(u_g\phi_h=\phi_h\).
Thus \(u_g\) is the identity. But \(g_D\iota_t=\iota_t g\);
naturality with the first factor gives
\[
 u_g\phi_t=\phi_t\,g_\Omega,
\]
where \(g_\Omega\) is the specialization of the prescribed
stabilizer automorphism. Consequently \(g_\Omega\) is the identity.

On the special fiber the pure stabilizer element \(g\) is its
specified automorphism of the Honda formal group over \(k_t\).
This is the special-fiber content of Lubin--Tate deformation
\cite[Proposition~3.3 and paragraph~3.4]{lubin-tate1966}, with its
coherent realization as above. If \(g\ne1\), some coefficient of
the power series \(g(X)-X\) is nonzero in \(k_t\). It remains
nonzero under the injective map \(k_t\to\Omega\), contrary to
\(g_\Omega=1\). Therefore \(g=1\).

This proof used only the coefficient map to \(\Omega\) for algebraic
base change; it did not require a ring-spectrum map to a spectrum
with coefficient field \(\Omega\). The assertion concerns equality
of actual field-valued maps; invariance of their kernels alone
would not give the conclusion needed next.
\end{proof}

\begin{lemma}[Algebraic trace one]\label{wup:algebraic-trace}
Let a finite group \(H\) act on a commutative ring \(R\). Suppose
that for every unital field-valued map \(x:R\to\Omega\),
\(x\circ g=x\) implies \(g=1\). Then
\[
 1\in\operatorname{im}\bigl(\Tr:R\to R^H\bigr),
 \qquad \Tr(r)=\sum_{g\in H}g(r).
\]
\end{lemma}

\begin{proof}
The trace image is an ideal of \(R^H\), since trace is additive
and \(a\Tr(r)=\Tr(ar)\) for \(a\in R^H\). Suppose this ideal
is proper, and choose a maximal ideal \(\mathfrak q\) containing
it. For every \(r\in R\), the monic orbit polynomial
\(\prod_{g\in H}(X-g(r))\) has coefficients in \(R^H\).
Thus \(R\) is integral over \(R^H\). Lying over gives a prime
\(\mathfrak p\) of \(R\) above \(\mathfrak q\). The quotient and
inclusion into its fraction field give a unital map
\(x:R\to\Omega=\operatorname{Frac}(R/\mathfrak p)\) killing
\(\mathfrak q\).

The multiplicative unital maps
\(\chi_g:R\to\Omega\), \(\chi_g(r)=x(g(r))\), are pairwise
distinct by the hypothesis. Such distinct maps are linearly
independent as functions. To prove this, suppose
\(\sum_{i=1}^v a_i\chi_i=0\) is a nonzero relation with the
least possible number of nonzero coefficients. One term is
impossible by evaluation at \(1\). If \(v>1\), choose an element
\(b\in R\) distinguishing \(\chi_1\) from \(\chi_v\).
Subtract \(\chi_v(b)\) times the relation at \(r\) from the
relation at \(br\). The result is
\[
 \sum_{i<v}a_i\bigl(\chi_i(b)-\chi_v(b)\bigr)\chi_i(r)=0.
\]
Its coefficient at \(i=1\) is nonzero in the field, giving a
shorter nonzero relation, a contradiction.

In particular the function \(\sum_{g\in H}\chi_g\) is not zero:
each of its coefficients is \(1\) in the field, even if the
characteristic divides \(|H|\). On the other hand it equals
\(x\circ\Tr\), which is zero since \(x\) kills the trace ideal.
This contradiction proves the claim. If \(R\) is the zero ring,
the conclusion holds directly with \(1=0\).
\end{proof}

\begin{proof}[Proof of Proposition~\ref{wup:completed-trace}]
Lemmas~\ref{wup:field-points} and \ref{wup:algebraic-trace}
give \(\bar z\in R\) with \(\Tr(\bar z)=1\). Lift it to
\(z_0\in\pi_0D\), and put \(v=\sum_g g(z_0)\). Then
\[
 v-1\in\m_t\pi_0D.
\]
The mixed module \(\pi_*D\) is flat over \((E_t)_*\) by
Lemma~\ref{wup:mixed-flat}. Lemma~\ref{wup:completion-reduction}
therefore identifies the homotopy of its actual localization
\(D\to D'\) with the ordinary \(\m_t\)-adic completion of
\(\pi_*D\).

At the finite adic stage \(\pi_0D/\m_t^j\pi_0D\), the element
\[
 w_j=\sum_{k=0}^{j-1}(1-v)^k
\]
satisfies \(v w_j=1-(1-v)^j=1\). These inverses are compatible
under reduction. Here \(j\) is a finite truncation level; the
quotient rings themselves are not required to be finite sets.
The comparison with completion is induced by the actual ring map
\(D\to D'\). Multiplication by the image of \(v\) is
\(O_t\)-linear and adically continuous and agrees with
multiplication by \(v\) on the dense ordinary image. It hence
agrees with the levelwise multiplication on the completion.
The compatible \(w_j\) give an inverse to the image of \(v\) in
\(\pi_0D'\).

Localization is equivariant, so that image of \(v\) is the
\(H\)-trace of the image of \(z_0\). Lemma~\ref{wup:trace-split}
now applies. If the completed ring is zero, its unit is zero and
every unital module is contractible, giving the same conclusion.
\end{proof}
